% RECYM — IEEE PES Transactions draft (IEEEtran)
% Target: IEEE Transactions on Power Delivery
% First submission limit (PES): <= 10 pages
\documentclass[journal]{IEEEtran}

\usepackage{cite}
\usepackage{amsmath,amssymb,amsfonts}
\usepackage{graphicx}
\usepackage{booktabs}
\usepackage{url}
% hyperref omitted for lean MiKTeX builds

\begin{document}

\title{RECYM: A Multi-Feeder CYMDIST Workflow for Demand Allocation,
Spot-Load Studies, and Situational--Projected Load-Flow Comparison}

\author{\IEEEauthorblockN{First~A.~Author\IEEEauthorrefmark{1},
Second~B.~Author\IEEEauthorrefmark{2}, and
Third~C.~Author\IEEEauthorrefmark{1}}
\IEEEauthorblockA{\IEEEauthorrefmark{1}TODO Affiliation 1, City, Country
(e-mail: corresponding@author.edu)}
\IEEEauthorblockA{\IEEEauthorrefmark{2}TODO Affiliation 2, City, Country}}

\maketitle

\begin{abstract}
Distribution utilities must repeatedly correct network models, allocate measured
headroom demand, connect prospective concentrated loads, and compare base versus
future load-flow cases inside commercial tools such as CYMDIST. Manual campaigns
are error-prone, hard to audit, and difficult to reproduce across dozens of
feeders. This paper presents RECYM, a multi-feeder software suite that couples
CYMDIST/CymPy and COM engines with a seven-step single-page application and a
batch pipeline. The workflow enforces a model-quality gate, maps large-customer
energy and contracted power onto secondary distribution transformers, allocates
residual demand by energy (kWh) against a measured headroom total, inserts locked
three-phase spot loads on existing nodes, and runs dual load-flow scenarios that
disconnect or reconnect the prospective load. On Electro Dunas feeder PA217
(22.9\,kV; 1324 nodes; 253 loads), RECYM preserved headroom balance at
9537.88\,kW, updated ten secondary transformers with energy-to-kWh and
power-to-kW locked injections, and connected a 1420\,kW spot load for
situational--projected comparison, with zero diagnostic problems on the quality
board after corrections. The architecture, coupling rules, and reproducible
artifacts are documented so peer utilities can replicate the campaign on
additional feeders.
\end{abstract}

\begin{IEEEkeywords}
CYMDIST, demand allocation, distribution networks, load flow,
power system modeling, software tools, spot load.
\end{IEEEkeywords}

\section{Introduction}
Planning and connection studies on medium-voltage distribution feeders require a
trustworthy digital twin: correct connectivity and equipment data, demand that
matches metered headroom, explicit treatment of large customers, and a transparent
comparison between the present operating point and a projected case that includes
a new concentrated load~\cite{kersting2017,gonen2014}.
Commercial packages such as CYMDIST provide load allocation and load-flow
engines, but production campaigns still rely on long sequences of manual clicks,
spreadsheets, and ad-hoc scripts.

Open platforms emphasize flexible simulation APIs~\cite{dugan2011}, while vendor
documentation covers allocation and load-flow algorithms in
isolation~\cite{cyme_alloc,cyme_lf}.
What remains scarce is an end-to-end operational workflow---diagnostics,
large-customer injection, kWh-based residual allocation, locked spot-load
insertion, dual-scenario load flow, and delivery-report generation---bound to a
multi-feeder configuration and an auditable user interface.

This paper presents RECYM. The contributions are:
(i)~a layered SPA/API/pipeline/CymPy--COM architecture for one feeder or batch
modes;
(ii)~a seven-step campaign with coupling rules that prevent incorrect mixing of
allocation, spot load, and situational/projected cases; and
(iii)~an industrial case study on feeder PA217 demonstrating headroom-consistent
allocation, SED updates, a 1420\,kW spot load, and reproducible artifacts.

\section{Related Work}
Distribution modeling fundamentals are well established~\cite{kersting2017,gonen2014}.
Load allocation to match metered feeder totals is standard utility
practice~\cite{cyme_alloc}.
Compensation-based power flow for weakly meshed networks underpins many
engines~\cite{shirmohammadi1988}.
OpenDSS~\cite{dugan2011} and pandapower~\cite{thurner2018} automate research
studies; RECYM instead orchestrates CYMDIST engines for industrial campaigns.

\section{System Architecture}
\subsection{Layered Design}
RECYM comprises: (1)~a React SPA with seven operator panels;
(2)~a FastAPI job layer with SSE progress and a Flask bridge;
(3)~domain pipelines (diagnostics, customers, allocation, spot load, load flow,
reports, optional optimization);
(4)~CymPy and COM adapters; and
(5)~external \texttt{.zxst}/\texttt{.mdb}/Excel data.
Fig.~\ref{fig:arch} (placeholder) shows the stack.

\subsection{Multi-Feeder Configuration}
Global paths live in \texttt{settings.json}; each feeder JSON stores
\texttt{network\_id}, study path, workbooks, and
\texttt{data/output/feeders/<ID>/}.
Writes require a resolved study path and \texttt{dry\_run=false}.

\section{Proposed Workflow}
\subsection{Seven-Step Campaign}
Table~\ref{tab:steps} summarizes the campaign.
After connecting a spot load, residual allocation must not be re-run; only load
flow follows.
Situational versus projected differs solely by disconnecting or reconnecting the
prospective spot load.

\begin{table}[!t]
\caption{RECYM Operator Campaign}
\label{tab:steps}
\centering
\begin{tabular}{@{}clp{5.6cm}@{}}
\toprule
Step & Panel & Function \\
\midrule
1 & Context & Bind DB/study; save headroom $P,Q$ \\
2 & Quality & NetworkDiagnostic; corrections; gate \\
3 & Customers & SED EA$\to$kWh, Pot$\to$kW; allocation \\
4 & Spot load & Existing node; $P/3$ per phase; Locked \\
5 & Load flows & Situational (off) / projected (on) \\
6 & Reports & Meta OCR + Word/PDF package \\
7 & Suite & Optional equipment optimization \\
\bottomrule
\end{tabular}
\end{table}

\subsection{Spot-Load Phase Split}
For three-phase demand $P$ and $Q$ at an existing node,
\begin{equation}
P_{\phi}=\frac{P}{3},\qquad Q_{\phi}=\frac{Q}{3},\quad \phi\in\{A,B,C\}.
\end{equation}

\section{Case Study Setup}
Electro Dunas feeder PA217 (\texttt{NET\_2030\_179\_PA217}),
$V_{\mathrm{ll}}=22.9$\,kV, has 1324 nodes, 1324 sections, and 253 loads.
Headroom: $P=9537.88$\,kW, $Q=2587.65$\,kvar, $S=9882.67$\,kVA
(loading factor 65.05\,\%).
Prospect load: 1420\,kW at node \texttt{NODE\_1080\_320357\_MI-17145\_2}.
Software stack: CYMDIST 9.2, CymPy, COM, RECYM SPA v6.

\section{Results and Discussion}
After corrections, the quality board reported zero problems.
Ten SED loads were written Connected with EA$\to$Consumo and Pot$\to$kW Locked;
one customer was excluded (Disconnected); one row lacked SED mapping.
Allocation validation: $\sum P=9537.87$\,kW versus headroom $9537.88$\,kW
($d=0.0$\,\%).
The 1420\,kW spot load uses equal phase split
($P_{\phi}=473.33$\,kW/phase).
Integral validation for PA217 returned PASS (21 September 2026).
Load-flow $V_{\min}$/loss scalars must be frozen from a non-dry-run export
before camera-ready tables.
Limitations include Windows COM dependence, single-session assumptions, and
utility-specific workbook mapping.

\section{Conclusion}
RECYM orchestrates CYMDIST for multi-feeder demand and connection studies via a
quality gate, SED injection, kWh allocation, locked spot loads, and
situational--projected load flows.
Future work includes multi-feeder LF KPI packs and production use of optional
equipment optimization.

\section*{Acknowledgment}
TODO: funding and institutional support.

\begin{thebibliography}{00}
\bibitem{kersting2017} W.~H. Kersting, \emph{Distribution System Modeling and Analysis}, 4th~ed. Boca Raton, FL, USA: CRC Press, 2017.
\bibitem{gonen2014} T.~G\"onen, \emph{Electric Power Distribution Engineering}, 3rd~ed. Boca Raton, FL, USA: CRC Press, 2014.
\bibitem{dugan2011} R.~C. Dugan and T.~E. McDermott, ``An open source platform for collaborating on smart grid research,'' in \emph{Proc. IEEE Power Energy Soc. Gen. Meeting}, Detroit, MI, USA, 2011, pp.~1--7.
\bibitem{cyme_alloc} CYME International T\&D, \emph{CYMDIST User Guide} (Load Allocation), Eaton, ver.~9.2.
\bibitem{cyme_lf} CYME International T\&D, \emph{CYMDIST User Guide} (Balanced Load Flow), Eaton, ver.~9.2.
\bibitem{shirmohammadi1988} D.~Shirmohammadi, H.~W. Hong, A.~Semlyen, and G.~X. Luo, ``A compensation-based power flow method for weakly meshed distribution and transmission networks,'' \emph{IEEE Trans. Power Syst.}, vol.~3, no.~2, pp.~753--762, May 1988.
\bibitem{thurner2018} L.~Thurner \emph{et al.}, ``pandapower---An open-source Python tool for convenient modeling, analysis, and optimization of electric power systems,'' \emph{IEEE Trans. Power Syst.}, vol.~33, no.~6, pp.~6510--6521, Nov. 2018.
\end{thebibliography}

\end{document}
